\frfilename{mt111.tex}
\versiondate{26.1.05}
\copyrightdate{1994}

\def\chaptername{Ruang ukur}
\def\sectionname{Aljabar-$\sigma$}

\newsection{111}

Dalam pengantar bab ini saya menyatakan bahwa ruang ukur adalah
`suatu himpunan yang beberapa (biasanya tidak semua) subhimpunannya
dapat diberi ukuran'. Semua konsep biasa tentang `panjang', `luas',
atau `volume' hanya berlaku pada himpunan yang cukup teratur. Teori
ukuran modern luar biasa kuat karena begitu beragam himpunan ternyata
cukup teratur untuk dapat diukur; tetapi kita tetap
harus mengharapkan adanya batasan, dan ketika mempelajari suatu ukuran,
pemahaman yang tepat tentang kelas himpunan yang diukurnya akan menjadi
pokok pekerjaan kita. Definisi dasarnya di sini ialah `aljabar-$\sigma$
atas himpunan'; semua ukuran dalam teori standar didefinisikan pada
koleksi semacam itu. Karena itu, saya mulai dengan menyatakan definisi
tersebut dan membahas secara singkat sifat-sifat kelas ini.


\leader{111A}{Definisi} Misalkan $X$ suatu himpunan. Sebuah
{\bf aljabar-$\sigma$ pada $X$}\cmmnt{ (kadang
disebut {\bf medan-$\sigma$})} adalah keluarga $\Sigma$ yang terdiri
atas subhimpunan-subhimpunan $X$ sedemikian sehingga

\qquad(i) $\emptyset\in\Sigma$;

\qquad(ii) untuk setiap $E\in\Sigma$, komplemennya $X\setminus E$ dalam
$X$ termasuk dalam $\Sigma$;

\qquad(iii) untuk setiap barisan $\langle E_n\rangle_{n\in\Bbb N}$ di
dalam $\Sigma$, gabungannya $\bigcup_{n\in\Bbb N}E_n$ termasuk dalam
$\Sigma$.


\cmmnt{
\leader{111B}{Catatan (a)} Hampir setiap bidang baru dalam matematika
murni kemungkinan besar diawali dengan definisi. Pada tahap ini, tidak
ada pengganti untuk menghafalkannya. Definisi-definisi ini merangkum
pemikiran banyak orang selama bertahun-tahun, kadang berabad-abad; Anda
tidak dapat mengharapkan semuanya selalu sesuai dengan gagasan yang
sudah akrab.


\header{111Bb}{\bf (b)} Meskipun demikian, Anda hendaknya selalu segera
mencari cara untuk membuat definisi baru menjadi lebih konkret dengan
menemukan contoh-contoh dari pengalaman matematika Anda sebelumnya.
Dalam hal `aljabar-$\sigma$', contoh-contoh yang benar-benar penting,
yang akan dijelaskan di bawah, pada dasarnya akan merupakan hal baru
--- andaikan, tentu saja, Anda memang perlu membaca bab ini. Namun, dua
contoh semestinya langsung dapat Anda pahami, dan mulai sekarang
hendaknya Anda selalu mengingatnya:

\qquad(i) untuk setiap $X$, $\Sigma=\{\emptyset,X\}$ adalah
aljabar-$\sigma$ pada $X$;

\qquad(ii) untuk setiap $X$, $\Cal PX$, yaitu himpunan semua subhimpunan
$X$, adalah aljabar-$\sigma$ pada $X$.

\noindent Tentu saja, keduanya masing-masing merupakan
aljabar-$\sigma$ terkecil dan terbesar atas subhimpunan-subhimpunan $X$.
Meskipun kita hanya akan mencurahkan sedikit waktu untuk keduanya,
sebenarnya keduanya sama-sama penting.


\header{111Bc}{\bf *(c)} Istilah {\bf ruang terukur} sering digunakan
untuk menyatakan pasangan $(X,\Sigma)$, dengan $X$ suatu himpunan dan
$\Sigma$ suatu aljabar-$\sigma$ pada $X$; tetapi
saya sendiri lebih suka menghindari istilah ini, kecuali ketika sangat
terdesak oleh waktu, sebab banyak di antara contoh paling menarik dari
objek semacam itu sebenarnya tidak mempunyai ukuran berguna yang
terkait dengannya.
}%end of comment

\cmmnt{
\leader{111C}{Gabungan dan irisan tak hingga} Jika Anda belum pernah
menjumpai gabungan tak hingga, ada baiknya berhenti sejenak untuk
mencermati rumus $\bigcup_{n\in\Bbb N}E_n$. Ini adalah himpunan titik
yang termasuk dalam satu atau lebih di antara himpunan-himpunan $E_n$;
kita dapat menuliskannya sebagai

$$\eqalign{\bigcup_{n\in\Bbb N}E_n
&=\{x:\exists\enskip n\in\Bbb N,\,x\in E_n\}\cr
&=E_0\cup E_1\cup E_2\cup\ldots.\cr}$$

\noindent (Saya menggunakan $\Bbb N$ untuk himpunan bilangan asli
$\{0,1,2,3,\ldots\}$.) Dengan cara yang sama,

$$\eqalign{\bigcap_{n\in\Bbb N}E_n
&=\{x:x\in E_n\Forall n\in\Bbb N\}\cr
&=E_0\cap E_1\cap E_2\cap\ldots.\cr}$$

\noindent Salah satu ciri teori dasar ruang ukur ialah bahwa teori ini
menuntut kemahiran lebih besar dalam operasi himpunan $\cup$, $\cap$,
$\setminus$ (`selisih himpunan': $E\setminus F=\{x:x\in
E,\,x\notin F\}$), $\symmdiff$ (`selisih simetris': $E\symmdiff
F=(E\setminus F)\cup(F\setminus E)=(E\cup F)\setminus(E\cap F)$)
daripada yang mungkin pernah Anda perlukan sebelumnya, ditambah lagi
dengan kerumitan gabungan dan irisan tak hingga. Saya sangat
menyarankan agar pada suatu saat Anda meluangkan setidaknya sedikit
waktu untuk Latihan 111Xa.
}%end of comment

\leader{111D}{Sifat-sifat dasar aljabar-$\sigma$} Jika $\Sigma$ adalah
aljabar-$\sigma$ pada $X$, maka aljabar tersebut
mempunyai sifat-sifat berikut.

\header{111Da}{\bf (a)} $E\cup F\in\Sigma$ untuk semua $E$,
$F\in\Sigma$.
\prooflet{\Prf\ Sebab, jika $E$, $F\in\Sigma$, tetapkan $E_0=E$,
$E_n=F$ untuk $n\ge 1$; maka $\langle E_n\rangle_{n\in\Bbb N}$ adalah
barisan di dalam $\Sigma$ dan $E\cup
F=\bigcup_{n\in\Bbb N}E_n\in\Sigma$.\ \Qed}


\header{111Db}{\bf (b)} $E\cap F\in\Sigma$ untuk semua $E$,
$F\in\Sigma$.
\prooflet{\Prf\ Menurut (ii) dalam definisi 111A, $X\setminus E$ dan
$X\setminus F\in\Sigma$; menurut (a) pada paragraf ini,
$(X\setminus E)\cup(X\setminus F)\in\Sigma$; kembali menurut 111A(ii),
$X\setminus((X\setminus E)\cup(X\setminus F))\in\Sigma$; tetapi
himpunan terakhir ini tidak lain adalah $E\cap F$.\ \Qed}


\header{111Dc}{\bf (c)} $E\setminus F\in\Sigma$ untuk semua $E$,
$F\in\Sigma$.
\prooflet{\Prf\ $E\setminus F=E\cap(X\setminus F)$.\ \Qed}


\header{111Dd}{\bf (d)} Sekarang, misalkan
$\langle E_n\rangle_{n\in\Bbb N}$ adalah barisan di dalam $\Sigma$, dan
perhatikan bahwa

$$\eqalign{\bigcap_{n\in\Bbb N}E_n
&=\{x:x\in E_n\Forall n\in\Bbb N\}\cr
&=E_0\cap E_1\cap E_2\cap\ldots\cr
&=X\setminus\bigcup_{n\in\Bbb N}(X\setminus E_n);\cr}$$

\noindent himpunan ini juga termasuk dalam $\Sigma$.


\cmmnt{
\leader{111E}{Lebih lanjut tentang gabungan dan irisan tak hingga (a)}
Sejauh ini, saya baru membahas gabungan dan irisan tak hingga dalam
konteks barisan $\langle E_n\rangle_{n\in\Bbb N}$ yang diindeks oleh
himpunan bilangan asli $\Bbb N$ itu sendiri. Banyak bentuk lain akan
muncul secara kurang lebih alami pada halaman-halaman berikutnya.
Sebagai contoh, perhatikan himpunan-himpunan berbentuk

\Centerline{$\bigcup_{n\ge 4}E_n=E_4\cup E_5\cup E_6\cup\ldots$,}

\Centerline{$\bigcup_{n\in\Bbb Z}E_n
=\{x:\exists\enskip n\in\Bbb Z,\,x\in E_n\}
=\ldots\cup E_{-2}\cup E_{-1}\cup E_0\cup E_1\cup E_2\cup\ldots$,}

\Centerline{$\bigcup_{q\in\Bbb Q}E_q
=\{x:\exists\enskip q\in\Bbb Q,\,x\in E_q\}$,}

\noindent di sini saya menggunakan $\Bbb Z$ untuk himpunan semua
bilangan bulat dan $\Bbb Q$ untuk himpunan bilangan rasional. Jika
setiap $E_n$, $E_q$ termasuk dalam suatu aljabar-$\sigma$ $\Sigma$,
gabungan-gabungan ini pun termasuk di dalamnya. Sebaliknya,

\Centerline{$\bigcup_{t\in[0,1]}E_t
=\{x:\exists\enskip t\in [0,1],\,x\in E_t\}$}

\noindent mungkin tidak termasuk dalam suatu aljabar-$\sigma$ yang
memuat setiap $E_t$. Karena itu, sangatlah penting untuk membangun
intuisi mengenai himpunan indeks yang `aman' dalam konteks ini, seperti
$\Bbb N$, $\Bbb Z$, dan $\Bbb Q$, serta himpunan indeks yang tidak aman.

\spheader 111Eb Saya sangat berharap bahwa Anda telah mempelajari teori
Cantor tentang himpunan tak hingga secukupnya sehingga catatan berikut
hanya mengungkapkan kembali hal-hal yang sudah akrab; tetapi jika belum,
saya berharap catatan ini dapat menjadi pengantar pertama, meskipun
sangat terbatas, untuk gagasan-gagasan tersebut. Inti dari tiga contoh
pertama ialah bahwa kita dapat mengindeks ulang keluarga-keluarga
himpunan yang terlibat sebagai barisan himpunan biasa. Untuk contoh
pertama, caranya sederhana; tuliskan $E'_n=E_{n+4}$ untuk
$n\in\Bbb N$, lalu perhatikan bahwa
$\bigcup_{n\ge 4}E_n=\bigcup_{n\in\Bbb N}E'_n\in\Sigma$.
Untuk dua contoh lainnya, kita perlu mengetahui sesuatu tentang
himpunan $\Bbb Z$ dan $\Bbb Q$. Kita dapat menemukan barisan bilangan
bulat $\langle k_n\rangle_{n\in\Bbb N}$ dan barisan bilangan rasional
$\langle q_n\rangle_{n\in\Bbb N}$ sedemikian sehingga setiap bilangan
bulat muncul (sekurang-kurangnya sekali) sebagai suatu $k_n$, dan setiap
bilangan rasional muncul (sekurang-kurangnya sekali) sebagai suatu
$q_n$; artinya, fungsi $n\mapsto k_n:\Bbb N\to\Bbb Z$ dan
$n\mapsto q_n:\Bbb N\to\Bbb Q$ bersifat surjektif.
\prooflet{\Prf\ Ada banyak cara untuk melakukannya; salah satunya ialah
menetapkan

$$\eqalign{k_n&=\Bover{n}2\text{ untuk }n\text{ genap},\cr
&=-\Bover{n+1}2\text{ untuk }n\text{ ganjil},\cr
q_n&=\Bover{n-m^3-m^2}{m+1}
  \text{ jika }m\in\Bbb N\text{ dan }m^3\le n<(m+1)^3.\cr}$$

\noindent (Anda sebaiknya memeriksa dengan saksama bahwa rumus-rumus ini
benar-benar bekerja seperti yang saya nyatakan.) \QeD}
Sekarang, untuk menangani $\bigcup_{n\in\Bbb Z}E_n$, kita dapat
menetapkan

\Centerline{$E'_n=E_{k_n}\in\Sigma$}

\noindent untuk $n\in\Bbb N$, sehingga

\Centerline{$\bigcup_{n\in\Bbb Z}E_n=\bigcup_{n\in\Bbb N}E_{k_n}
=\bigcup_{n\in\Bbb N}E'_n\in\Sigma$,}

\noindent sedangkan untuk kasus yang lain kita mempunyai

\Centerline{$\bigcup_{q\in\Bbb Q}E_q
=\bigcup_{n\in\Bbb N}E_{q_n}\in\Sigma$.}

Perhatikan bahwa kasus pertama $\bigcup_{n\ge 4}E_n$ dapat dipandang
sebagai penerapan asas yang sama; pemetaan $n\mapsto n+4$ adalah
surjeksi dari $\Bbb N$ ke $\{4,5,6,7,\ldots\}$.
}%end of comment

\vleader{48pt}{111F}{Himpunan terhitung (a)}\cmmnt{ Ciri yang sama pada
himpunan-himpunan $\{n:n\ge 4\}$, $\Bbb Z$, dan $\Bbb Q$ yang
memungkinkan prosedur ini ialah bahwa semuanya `terhitung'. Untuk
keperluan kita di sini, definisi keterhitungan yang paling alami adalah
sebagai berikut:} \dvro{S}{s}uatu himpunan $K$ disebut {\bf terhitung}
jika himpunan tersebut kosong atau terdapat suatu surjeksi dari $\Bbb N$ ke $K$. Dalam
hal ini, jika $\Sigma$ adalah suatu aljabar-$\sigma$
dan $\langle E_k\rangle_{k\in K}$ adalah keluarga di dalam $\Sigma$
yang diindeks oleh $K$, maka $\bigcup_{k\in K}E_k\in\Sigma$.
\prooflet{\Prf\ Sebab, jika $n\mapsto k_n:\Bbb N\to K$ adalah suatu
surjeksi, maka $E'_n=E_{k_n}\in\Sigma$ untuk setiap $n\in\Bbb N$, dan
$\bigcup_{k\in K}E_k=\bigcup_{n\in\Bbb N}E'_n\in\Sigma$. Ini belum
mencakup kasus $K=\emptyset$; tetapi dalam kasus ini, penafsiran alami
untuk $\bigcup_{k\in K}E_k$ ialah

\Centerline{$\{x:\exists\enskip k\in \emptyset,\,x\in E_k\}$}

\noindent yang juga merupakan $\emptyset$, dan karenanya termasuk dalam
$\Sigma$ menurut syarat (i) dari 111A.\ \QeD} \cmmnt{(Dalam arti
tertentu, perlakuan terhadap $\emptyset$ ini merupakan masalah
kesepakatan; tetapi ada berbagai konteks tempat kita ingin membahas
$\bigcup_{k\in K}E_k$ tanpa memeriksa apakah $K$ benar-benar mempunyai
anggota, sehingga kita perlu memperjelas apa yang akan kita lakukan
dalam kasus seperti itu.)}

\header{111Fb}{\bf (b)}\cmmnt{ Terdapat teori yang luas dan sangat
penting mengenai himpunan terhitung. Bagian-bagian yang menurut saya
harus dinyatakan secara eksplisit pada tahap ini hanyalah yang berikut.
(Dalam \S1A1 saya menambahkan beberapa kata untuk menghubungkan
penyajian ini dengan pendekatan-pendekatan yang lazim.)

\medskip

\quad}{\bf (i)} Jika $K$ terhitung dan $L\subseteq K$, maka $L$
terhitung. \prooflet{\Prf\ Jika $L=\emptyset$, hal ini langsung jelas.
Jika tidak, ambil sembarang $l^*\in L$ dan suatu surjeksi
$n\mapsto k_n:\Bbb N\to K$ (tentu saja $K$ juga tidak kosong, sebab
$l^*\in K$); tetapkan $l_n=k_n$ jika $k_n\in L$, dan $l^*$ jika tidak;
maka $n\mapsto l_n:\Bbb N\to L$ adalah suatu surjeksi.\ \Qed}

\medskip

\quad{\bf (ii)} Produk Kartesius $\Bbb N\times\Bbb N=\{(m,n):m$,
$n\in\Bbb N\}$ terhitung. \prooflet{\Prf\ Untuk setiap $n\in\Bbb N$,
misalkan $k_n$, $l_n\in\Bbb N$ sedemikian sehingga
$n+1=2^{k_n}(2l_n+1)$; artinya, $k_n$ adalah pangkat $2$ dalam
faktorisasi prima $n+1$, dan $2l_n+1$ adalah bilangan (yang pasti ganjil)
$(n+1)/2^{k_n}$. Sekarang, $n\mapsto(k_n,l_n)$ adalah suatu surjeksi
dari $\Bbb N$ ke $\Bbb N\times\Bbb N$.\ \Qed}\
\cmmnt{Kelak penting bagi kita untuk mengetahui bahwa
$n\mapsto(k_n,l_n)$ sebenarnya merupakan bijeksi, sebagaimana mudah
diperiksa.}

\medskip

\quad{\bf (iii)} Akibatnya, jika $K$ dan $L$ adalah himpunan terhitung,
maka $K\times L$ juga terhitung. \prooflet{\Prf\ Jika $K$ atau $L$
kosong, maka $K\times L$ kosong, sehingga dalam kasus ini $K\times L$
tentu terhitung. Jika tidak, misalkan $\phi:\Bbb N\to K$ dan
$\psi:\Bbb N\to L$ adalah surjeksi; maka kita mempunyai surjeksi
$\theta:\Bbb N\times\Bbb N\to K\times L$ yang didefinisikan dengan
menetapkan $\theta(m,n)=(\phi(m),\psi(n))$ untuk semua $m$,
$n\in\Bbb N$. Dari (ii) tepat di atas, kita tahu bahwa ada pula
surjeksi $\chi:\Bbb N\to\Bbb N\times\Bbb N$, sehingga
$\theta\chi:\Bbb N\to K\times L$ adalah suatu surjeksi, dan
$K\times L$ pasti terhitung.\ \Qed}

\medskip

\quad{\bf (iv)}\dvro{ Jika}{ Sekarang, induksi pada $r$ menunjukkan
bahwa jika} $K_1$, $K_2,\ldots,K_r$ adalah himpunan-himpunan terhitung,
maka $K_1\times\ldots\times K_r$ juga terhitung. Secara khusus,\cmmnt{
himpunan-himpunan seperti} $\Bbb Q^r\times\Bbb Q^r$ akan terhitung,
untuk setiap bilangan bulat $r\ge 1$.


\header{111Fc}{\bf (c)}\dvro{ Jika}{ Dengan menggabungkan 111Dd di atas
dengan gagasan-gagasan ini, kita melihat bahwa jika} $\Sigma$ adalah
suatu aljabar-$\sigma$, $K$ adalah himpunan terhitung
yang tidak kosong, dan $\langle E_k\rangle_{k\in K}$ adalah keluarga di
dalam $\Sigma$, maka

\Centerline{$\bigcap_{k\in K}E_k=\{x:x\in E_k\Forall k\in
K\}$}

\noindent termasuk dalam $\Sigma$. \prooflet{\Prf\ Misalkan
$n\mapsto k_n:\Bbb N\to K$ adalah suatu surjeksi; maka
$\bigcap_{k\in K}E_k=\bigcap_{n\in\Bbb N}E_{k_n}\in\Sigma$, seperti
pada 111Dd.\ \Qed}

\cmmnt{Perhatikan bahwa terdapat kesulitan pada pengertian
$\bigcap_{k\in K}E_k$ jika $K=\emptyset$; penafsiran alami rumus ini
ialah membacanya sebagai kelas semesta. Jadi, biasanya, jika ada
kemungkinan bahwa $K$ kosong, kita memerlukan perumusan seperti
$X\cap\bigcap_{k\in K}E_k$.
}%end of comment

\cmmnt{\header{111Fd}{\bf (d)} Sebagai contoh cara menggunakan
gagasan-gagasan ini, perhatikan hal berikut. Misalkan $X$ suatu
himpunan, $\Sigma$ suatu aljabar-$\sigma$ pada $X$, dan
$\langle E_{qn}\rangle_{q\in\Bbb Q,n\in\Bbb N}$ suatu
keluarga di dalam $\Sigma$. Maka

\Centerline{$E
=\bigcap_{q\in\Bbb Q,q<\sqrt{2}}\bigcup_{m\in\Bbb N}\bigcap_{n\ge m}E_{qn}
=\bigcap_{q\in\Bbb Q,q<\sqrt{2}}(\bigcup_{m\in\Bbb N}(\bigcap_{n\ge m}E_{qn}))
\in\Sigma$.}

\noindent\vthsp\prooflet{\Prf\ Tetapkan
$F_{qm}=\bigcap_{n\ge m}E_{qn}=\bigcap_{n\in\Bbb N}E_{q,m+n}$ untuk
$q\in\Bbb Q$ dan $m\in\Bbb N$; maka setiap $F_{qm}$ termasuk dalam
$\Sigma$, menurut 111Dd atau (c) di atas. Tetapkan
$G_q=\bigcup_{m\in\Bbb N}F_{qm}$ untuk $q\in\Bbb Q$; maka setiap $G_q$
termasuk dalam $\Sigma$, menurut 111A(iii). Tetapkan
$K=\{q:q\in\Bbb Q,\,q<\sqrt{2}\}$; maka $K$ terhitung, menurut 111E dan
(b-i) pada paragraf ini. Jadi, $\bigcap_{q\in K}G_q$ termasuk dalam
$\Sigma$, menurut (c). Tetapi $E=\bigcap_{q\in K}G_q$.\ \Qed}
}%end of comment

\cmmnt{\header{111Fe}{\bf (e)} Dan satu catatan terakhir, yang saya
berikan di sini tanpa bukti --- walaupun banyak bukti akan tersirat
dalam pembahasan di bawah, dan satu di antaranya saya uraikan dalam
1A1Ha ---

\medskip

\centerline{\bf Himpunan $\Bbb R$ dari semua bilangan real tidak
terhitung.}

\medskip

\noindent Jadi, Anda harus menahan godaan untuk mencari suatu daftar
$a_0$, $a_1,\ldots$ yang mencakup seluruh himpunan bilangan real.
}%end of comment

\leader{111G}{Himpunan Borel }\cmmnt{Di sini saya dapat menjelaskan
salah satu jenis aljabar-$\sigma$ taktrivial; perumusannya cukup
abstrak, tetapi tekniknya penting dan istilahnya merupakan bagian dari
kosakata dasar teori ukuran.

\medskip

}{\bf (a)} Misalkan $X$ suatu himpunan, dan misalkan $\frak S$ sembarang
keluarga tak kosong yang terdiri atas aljabar-$\sigma$ pada $X$.
\cmmnt{(Jadi, suatu {\it anggota}
$\frak S$ sendiri merupakan suatu {\it keluarga} himpunan;
$\frak S\subseteq\Cal P(\Cal PX)$.)} Maka

\Centerline{$\bigcap\frak S=\{E:E\in\Sigma$ untuk setiap $\Sigma\in\frak
S\}$,}

\noindent yaitu irisan semua aljabar-$\sigma$ yang termasuk dalam
$\frak S$, merupakan aljabar-$\sigma$ pada $X$.
\prooflet{\Prf\
(i) Menurut hipotesis, $\frak S$ tidak kosong; ambil
$\Sigma_0\in\frak S$; maka
$\bigcap\frak S\subseteq\Sigma_0\subseteq\Cal PX$, sehingga setiap
anggota $\bigcap\frak S$ merupakan subhimpunan $X$. (ii)
$\emptyset\in\Sigma$ untuk setiap $\Sigma\in\frak S$, sehingga
$\emptyset\in\bigcap\frak S$. (iii) Jika $E\in\bigcap\frak S$, maka
$E\in\Sigma$ untuk setiap $\Sigma\in\frak S$, sehingga
$X\setminus E\in\Sigma$ untuk setiap $\Sigma\in\frak S$ dan
$X\setminus E\in\bigcap\frak S$. (iv) Misalkan
$\langle E_n\rangle_{n\in\Bbb N}$ sembarang barisan di dalam
$\bigcap\frak S$. Maka, untuk setiap $\Sigma\in\frak S$,
$\langle E_n\rangle_{n\in\Bbb N}$ adalah barisan di dalam $\Sigma$,
sehingga $\bigcup_{n\in\Bbb N}E_n\in\Sigma$; karena $\Sigma$ dipilih sembarang,
$\bigcup_{n\in\Bbb N}E_n\in\bigcap\frak S$.\ \Qed}

\header{111Gb}{\bf (b)} Sekarang, misalkan $\Cal A$ sembarang keluarga
subhimpunan $X$. Perhatikan

\Centerline{$\frak S=\{\Sigma:\Sigma$ adalah suatu aljabar-$\sigma$ atas
subhimpunan-subhimpunan $X$, $\Cal A\subseteq\Sigma\}$.}

\noindent\cmmnt{Menurut definisi, $\frak S$ adalah suatu keluarga yang
terdiri atas aljabar-$\sigma$ pada $X$; selain
itu, keluarga tersebut tidak kosong, karena $\Cal PX\in\frak S$. Jadi}
$\Sigma_{\Cal A}=\bigcap\frak S$ adalah aljabar-$\sigma$ pada
$X$\cmmnt{. Karena $\Cal A\subseteq\Sigma$
untuk setiap $\Sigma\in\frak S$, maka
$\Cal A\subseteq\Sigma_{\Cal A}$; dengan demikian,
$\Sigma_{\Cal A}$ sendiri termasuk dalam $\frak S$}; ini adalah
aljabar-$\sigma$ terkecil atas subhimpunan-subhimpunan $X$ yang memuat
$\Cal A$.

Kita mengatakan bahwa $\Sigma_{\Cal A}$ adalah aljabar-$\sigma$ pada
$X$ yang {\bf dibangkitkan oleh} $\Cal A$.

\medskip

\noindent{\bf Contoh (i)} Untuk setiap $X$, aljabar-$\sigma$ pada $X$
yang dibangkitkan oleh $\emptyset$ adalah
$\{\emptyset,X\}$.

\medskip

\quad{\bf (ii)} Aljabar-$\sigma$ pada $\Bbb N$
yang dibangkitkan oleh $\{\{n\}:n\in\Bbb N\}$ adalah $\Cal P\Bbb N$.

\header{111Gc}{\bf (c)(i)} Kita mengatakan bahwa suatu himpunan
$G\subseteq\Bbb R$ {\bf terbuka} jika untuk setiap $x\in G$ terdapat
$\delta>0$ sedemikian sehingga interval terbuka
$\ooint{x-\delta,x+\delta}$ termuat dalam $G$.

\medskip

\quad{\bf (ii)} Demikian pula, untuk setiap $r\ge 1$, kita mengatakan
bahwa suatu himpunan $G\subseteq\BbbR^r$ {\bf terbuka} dalam
$\BbbR^r$ jika untuk setiap $x\in G$ terdapat $\delta>0$ sedemikian
sehingga $\{y:\|y-x\|<\delta\}\subseteq G$, dengan, untuk
$z=(\zeta_1,\ldots,\zeta_r)\in\BbbR^r$, saya menuliskan
$\|z\|=\sqrt{\sum_{i=1}^r|\zeta_i|^2}$\cmmnt{; jadi $\|y-x\|$ tidak
lain adalah jarak Euklides biasa dari $y$ ke $x$}.

\header{111Gd}{\bf (d)} Sekarang, {\bf himpunan-himpunan Borel} dalam
$\Bbb R$, atau dalam $\BbbR^r$, tidak lain adalah anggota
aljabar-$\sigma$ pada $\Bbb R$ atau $\BbbR^r$ yang dibangkitkan oleh
keluarga himpunan terbuka dalam
$\Bbb R$ atau $\BbbR^r$; aljabar-$\sigma$ itu sendiri dalam masing-masing
kasus disebut {\bf aljabar-$\sigma$ Borel}.

\cmmnt{\header{111Ge}{\bf (e)} Sebagian pembaca akan merasa, dengan
tepat, bahwa pengembangan di sini hanya memberi sedikit gambaran
tentang seperti apa `sebenarnya' suatu himpunan Borel. (Himpunan terbuka
jauh lebih mudah; lihat 111Ye.) Sebenarnya, pentingnya konsep ini
sebagian besar berasal dari adanya cara-cara alternatif yang lebih
eksplisit dan, dalam suatu arti, lebih konkret untuk menjelaskan
himpunan Borel. Saya akan kembali ke topik ini dalam Bab 42 di Jilid 4.
}%end of comment

\exercises{
\leader{111X}{Latihan dasar $\pmb{>}$(a)} Berlatihlah menggunakan
aljabar gabungan dan irisan tak hingga sampai Anda dapat dengan yakin
menafsirkan rumus-rumus seperti

\Centerline{$E\cap(\bigcup_{n\in\Bbb N}F_n)$,
  \qquad$\bigcup_{n\in\Bbb N}(E_n\setminus F)$,
  \qquad$E\cup(\bigcap_{n\in\Bbb N}F_n)$,}

\Centerline{$\bigcup_{n\in\Bbb N}(E\setminus F_n)$,
  \qquad$E\setminus(\bigcup_{n\in\Bbb N}F_n)$,
  \qquad$\bigcap_{n\in\Bbb N}(E_n\setminus F)$,}

\Centerline{$E\setminus(\bigcap_{n\in\Bbb N}F_n)$,
  \qquad$\bigcap_{n\in\Bbb N}(E\cup F_n)$,
  \qquad$(\bigcup_{n\in\Bbb N}E_n)\setminus F$,}

\Centerline{$\bigcup_{n\in\Bbb N}(E\cap F_n)$,
  \qquad$(\bigcap_{n\in\Bbb N}E_n)\setminus F$,
  \qquad$\bigcap_{n\in\Bbb N}(E\setminus F_n)$,}

\Centerline{$(\bigcup_{n\in\Bbb N}E_n)\cap(\bigcup_{n\in\Bbb N}F_n)$,
  \qquad$\bigcap_{m,n\in\Bbb N}(E_m\setminus F_n)$,
  \qquad$(\bigcap_{n\in\Bbb N}E_n)\cup(\bigcap_{n\in\Bbb N}F_n)$,}

\Centerline{$\bigcap_{m,n\in\Bbb N}(E_m\cup F_n)$,
  \qquad$(\bigcap_{n\in\Bbb N}E_n)\setminus(\bigcup_{n\in\Bbb N}F_n)$,
  \qquad$\bigcup_{m,n\in\Bbb N}(E_m\cap F_n)$,}

\noindent dan, khususnya, dapat mengenali sembilan pasangan alami yang
dibentuk oleh rumus-rumus tersebut.

\sqheader 111Xb Dalam $\Bbb R$, tunjukkan bahwa semua `interval
terbuka' $\ooint{a,b}$, $\ooint{-\infty,b}$, $\ooint{a,\infty}$ adalah
himpunan terbuka, dan bahwa semua interval (terbatas ataupun tak
terbatas, terbuka, tertutup, ataupun setengah terbuka) merupakan
himpunan Borel.

\sqheader 111Xc Misalkan $X$ dan $Y$ adalah himpunan dan $\Sigma$
suatu aljabar-$\sigma$ pada $X$. Misalkan
$\phi:X\to Y$ suatu fungsi. Tunjukkan bahwa
$\{F:F\subseteq Y,\,\phi^{-1}[F]\in\Sigma\}$ adalah aljabar-$\sigma$
pada $Y$. (Lihat 1A1B untuk notasi yang
digunakan di sini.)

\sqheader 111Xd Misalkan $X$ dan $Y$ adalah himpunan dan $\Tau$ suatu
aljabar-$\sigma$ pada $Y$. Misalkan
$\phi:X\to Y$ suatu fungsi. Tunjukkan bahwa
$\{\phi^{-1}[F]:F\in\Tau\}$ adalah aljabar-$\sigma$ pada $X$.

\spheader 111Xe Misalkan $X$ suatu himpunan, $\Cal A$ suatu keluarga
subhimpunan $X$, dan $\Sigma$ aljabar-$\sigma$ pada $X$ yang
dibangkitkan oleh $\Cal A$. Misalkan
$Y$ suatu himpunan lain dan $\phi:Y\to X$ suatu fungsi. Tunjukkan bahwa
$\{\phi^{-1}[E]:E\in\Sigma\}$ adalah aljabar-$\sigma$ pada $Y$ yang
dibangkitkan oleh
$\{\phi^{-1}[A]:A\in\Cal A\}$.

\spheader 111Xf Misalkan $X$ suatu himpunan, $\Cal A$ suatu keluarga
subhimpunan $X$, dan $\Sigma$ aljabar-$\sigma$ pada $X$ yang
dibangkitkan oleh $\Cal A$. Misalkan
$Y\subseteq X$. Tunjukkan bahwa $\{E\cap Y:E\in\Sigma\}$ adalah
aljabar-$\sigma$ pada $Y$ yang dibangkitkan oleh
$\{A\cap Y:A\in\Cal A\}$.

\leader{111Y}{Latihan lanjutan (a)} Dalam $\BbbR^r$, dengan $r\ge 1$,
tunjukkan bahwa $G+a=\{x+a:x\in G\}$ terbuka setiap kali
$G\subseteq\BbbR^r$ terbuka dan $a\in\BbbR^r$. Selanjutnya, tunjukkan
bahwa $E+a$ adalah himpunan Borel setiap kali $E\subseteq\BbbR^r$
adalah himpunan Borel dan $a\in\BbbR^r$. \Hint{tunjukkan bahwa
$\{E:E+a\text{ adalah himpunan Borel}\}$ adalah aljabar-$\sigma$ yang memuat
semua himpunan terbuka.}

\header{111Yb}{\bf (b)} Misalkan $X$ suatu himpunan, $\Sigma$ suatu
aljabar-$\sigma$ pada $X$, dan $A$ sembarang
subhimpunan $X$. Tunjukkan bahwa $\{(E\cap
A)\cup(F\setminus A):E,\,F\in\Sigma\}$ adalah aljabar-$\sigma$ pada
$X$, yaitu aljabar-$\sigma$ yang dibangkitkan
oleh $\Sigma\cup\{A\}$.

\header{111Yc}{\bf (c)} Misalkan $G\subseteq\BbbR^2$ suatu himpunan
terbuka. Tunjukkan bahwa semua penampang horizontal dan vertikal

\Centerline{$\{\xi:(\xi,\eta)\in G\}$,
\quad$\{\xi:(\eta,\xi)\in G\}$}

\noindent dari $G$ merupakan subhimpunan terbuka dari $\Bbb R$.

\header{111Yd}{\bf (d)} Misalkan $E\subseteq\BbbR^2$ suatu himpunan
Borel. Tunjukkan bahwa semua penampang horizontal dan vertikal

\Centerline{$\{\xi:(\xi,\eta)\in E\}$,
\quad$\{\xi:(\eta,\xi)\in E\}$}

\noindent dari $E$ merupakan subhimpunan Borel dari $\Bbb R$.
\Hint{tunjukkan bahwa keluarga subhimpunan-subhimpunan $\BbbR^2$ yang
semua penampangnya merupakan himpunan Borel adalah aljabar-$\sigma$ pada
$\BbbR^2$ yang memuat semua himpunan terbuka.}

\spheader 111Ye Misalkan $G\subseteq\Bbb R$ suatu himpunan terbuka.
Tunjukkan bahwa $G$ dapat dinyatakan secara unik sebagai gabungan suatu
keluarga terhitung $\Cal I$ (yang mungkin kosong) yang terdiri atas
interval-interval terbuka (`komponen-komponen' $G$), dengan tidak ada
dua interval yang mempunyai titik bersama. \Hint{untuk $x$, $y\in G$,
nyatakan $x\sim y$ jika setiap titik di antara $x$ dan $y$ termasuk
dalam $G$. Tunjukkan bahwa $\sim$ merupakan relasi ekuivalensi.
Misalkan $\Cal I$ adalah himpunan kelas-kelas ekuivalensinya.}
}%end of exercises

\endnotes{
\Notesheader{111} Menurut saya, konsep terpenting dalam bagian ini
adalah konsep yang diperkenalkan sambil lalu dalam 111E--111F,
yaitu gagasan tentang himpunan `terhitung'. Walaupun untuk sementara
sebagian besar teori formal himpunan tak hingga dapat dihindari, menurut
saya mustahil memahami bab ini tanpa pemahaman yang kokoh mengenai
perbedaan antara `hingga' dan `tak hingga', serta intuisi tertentu
mengenai `keterhitungan'. Secara khusus, Anda harus mengingat bahwa
himpunan tak hingga pada umumnya tidak terhitung, dan bahwa
aljabar-$\sigma$ pada umumnya tidak tertutup terhadap gabungan sembarang.

Hal berikutnya yang perlu Anda pastikan ialah bahwa Anda dapat menangani
manipulasi teori himpunan di sini, sehingga rumus seperti
$\bigcap_{n\in\Bbb N}E_n=X\setminus\bigcup_{n\in\Bbb N}(X\setminus E_n)$
(111Dd), meskipun belum transparan, setidaknya tidak lagi terasa
menggentarkan. Sebagian besar isi volume ini akan dinyatakan dalam
bahasa tersebut, sehingga kefasihan yang memadai sangatlah penting.

Terakhir, bagi mereka yang mencari sebuah gagasan nyata untuk segera
digarap, saya menawarkan konsep aljabar-$\sigma$ yang `dibangkitkan'
oleh suatu koleksi $\Cal A$ (111G). Inti definisi di sini ialah bahwa
definisi tersebut melibatkan pembahasan suatu keluarga
$\frak S\in\Cal P(\Cal P(\Cal PX))$, meskipun $\Cal A$ dan
$\Sigma_{\Cal A}$ keduanya merupakan subhimpunan $\Cal PX$; kita perlu
bekerja satu atau dua lapis lebih tinggi dalam hierarki himpunan kuasa.
Sebagai contoh, Anda mungkin pernah menjumpai konsep `subruang linear
$U$ yang dibangkitkan oleh vektor-vektor $u_1,\ldots,u_n$'. Konsep ini
dapat didefinisikan sebagai irisan semua subruang linear yang memuat
vektor-vektor $u_1,\ldots,u_n$, yaitu metode yang bersesuaian dengan
metode 111Ga--b; tetapi konsep tersebut juga mempunyai definisi
`internal', yakni sebagai himpunan vektor yang dapat dinyatakan dalam
bentuk $\alpha_1u_1+\ldots+\alpha_nu_n$ untuk skalar-skalar $\alpha_i$.
Namun, untuk aljabar-$\sigma$, tidak tersedia definisi `internal'
sesederhana itu (walaupun ada banyak hal yang dapat dikatakan mengenai
arah ini yang menurut saya belum siap kita bahas; beberapa gagasan
dapat ditemukan dalam \S136). Penyebab utamanya ialah (iii) dalam
definisi 111A; suatu aljabar-$\sigma$ harus tertutup terhadap suatu
operasi infinitari, yakni operasi gabungan yang diterapkan pada barisan
himpunan tak hingga. Sebaliknya, subruang linear dari suatu ruang vektor
hanya perlu tertutup terhadap operasi finitari perkalian skalar dan
penjumlahan, yang masing-masing melibatkan paling banyak dua vektor pada
suatu waktu.
}%end of notes

\discrpage
